\documentclass[a4paper]{article}

\usepackage[fontset=fandol]{ctex}
\usepackage{amsmath, amssymb}
\usepackage{graphicx}
\usepackage[margin=2.2cm]{geometry}
\usepackage[most]{tcolorbox}
\usepackage{etoolbox}
\usepackage{listings}
\usepackage{hyperref}
\usepackage{xurl}
\usepackage{booktabs}
\usepackage{longtable}
\usepackage{tabularx}
\usepackage{array}
\usepackage{subcaption}
\usepackage{float}
\usepackage{enumitem}
\usepackage{tikz}
\usetikzlibrary{arrows.meta, positioning, shapes.geometric}

\hypersetup{
    colorlinks=true,
    linkcolor=blue!50!black,
    urlcolor=blue!60!black,
    citecolor=blue!50!black
}

\newtcolorbox{findingbox}[1]{
    enhanced,
    colback=green!5!white,
    colframe=green!45!black,
    colbacktitle=green!45!black,
    coltitle=white,
    fonttitle=\bfseries,
    title=#1,
    sharp corners
}

\newtcolorbox{evidencebox}[1]{
    enhanced,
    colback=blue!4!white,
    colframe=blue!65!black,
    colbacktitle=blue!65!black,
    coltitle=white,
    fonttitle=\bfseries,
    title=#1,
    sharp corners
}

\newtcolorbox{riskbox}[1]{
    enhanced,
    colback=red!4!white,
    colframe=red!70!black,
    colbacktitle=red!70!black,
    coltitle=white,
    fonttitle=\bfseries,
    title=#1,
    sharp corners
}

\newtcolorbox{methodbox}[1]{
    enhanced,
    colback=yellow!9!white,
    colframe=yellow!60!black,
    colbacktitle=yellow!60!black,
    coltitle=black,
    fonttitle=\bfseries,
    title=#1,
    sharp corners
}

\lstset{
    basicstyle=\ttfamily\small,
    breaklines=true,
    frame=single,
    numbers=left,
    numberstyle=\tiny\color{gray},
    captionpos=b,
    extendedchars=false
}

\newcolumntype{Y}{>{\raggedright\arraybackslash}X}
\setlist[itemize]{leftmargin=1.8em}
\newcommand{\filepath}[1]{\nolinkurl{#1}}

\newcommand{\reporttitle}{CCH 运行时网络故障与部署加固报告}
\newcommand{\reportsubtitle}{错误 DNS 覆盖、入口/出口链路混淆，以及跨平台 restart/probe 修复}
\newcommand{\reportauthors}{Codex}
\newcommand{\reportdate}{2026-03-25}
\newcommand{\researchquestion}{这次 RC/Hiyo 直连异常究竟坏在哪一层；为什么 Tailscale Serve 正常并不能推出容器可以出公网；以后如何更快定位和恢复？}
\newcommand{\reportscope}{运行时目录 \filepath{/home/fanghaotian/Applications/claude-code-hub}，仓库 \filepath{/home/fanghaotian/Desktop/src/cch/claude-code-hub}，证据窗口截至 2026-03-25 20:55 CST。}
\newcommand{\reportaudience}{claude-code-hub 维护者、部署脚本维护者、日常运维排障者}
\newcommand{\sourcewindow}{本地运行配置、容器日志、健康检查、连通性探测，以及仓库当前工作树，访问日期均为 2026-03-25}
\newcommand{\reportcoverpath}{}

\begin{document}

\begin{titlepage}
\centering
{\Large 深度研究报告\par}
\vspace{1.0cm}
{\huge\bfseries \reporttitle\par}
\vspace{0.5cm}
{\Large \reportsubtitle\par}
\vspace{0.9cm}
{\large \reportauthors\par}
\vspace{0.2cm}
{\large \reportdate\par}
\vspace{1.0cm}

{\small 本报告无外部封面图，主体证据来自运行时配置、容器日志和本地命令输出。\par}

\vfill
\begin{tcolorbox}[width=0.94\textwidth, colback=black!2!white, colframe=black!60, sharp corners]
\textbf{核心问题}：\researchquestion\par
\textbf{研究范围}：\reportscope\par
\textbf{目标读者}：\reportaudience\par
\textbf{证据窗口}：\sourcewindow\par
\end{tcolorbox}
\end{titlepage}

\tableofcontents
\newpage

\section{执行摘要}

\begin{findingbox}{核心结论}
这次 ``RC/Hiyo 不能直连'' 的主因并不是 provider 的 \texttt{base URL} 或选择逻辑被改坏，而是 \textbf{运行时 compose 给容器加了错误的硬编码 DNS}。保留下来的最强证据是备份文件 \filepath{/home/fanghaotian/Applications/claude-code-hub/docker-compose.archbox.yaml.bak-20260325-1912} 中的 \texttt{dns: 119.29.29.29 / 223.5.5.5}。这会把问题从 ``应用配置层'' 推到 ``容器出公网与解析链路层''。
\end{findingbox}

围绕这次事故，可以得到四个比 ``把配置改爆了'' 更精确的判断：

\begin{itemize}
\item \textbf{provider 选择本身没有坏。} 日志里两次成功请求都明确显示选中的仍然是 \texttt{RC}，而且 endpoint 仍然是 \texttt{https://right.codes/codex/}。
\item \textbf{Tailscale Serve 正常不等于容器能出公网。} 前者证明的是别的机器可以通过 tailnet 打到宿主机上的服务入口；后者要求容器能主动完成 DNS 解析、TCP/TLS 建连并访问外部站点，二者是不同方向的链路。
\item \textbf{当前已修复的部分是容器网络回到正常形态。} 现在 live 的 \texttt{claude-code-hub-app-s6gz} 没有 Docker 级别 \texttt{Dns} 覆盖，健康检查正常，\texttt{right.codes} 从 host 和容器都能访问。
\item \textbf{当前剩余的 Hiyo 症状并不等价于原事故。} \texttt{codex.hiyo.top} 现在在 host 和容器里都会报同样的 TLS EOF，这更像 endpoint 侧问题，而不是 CCH 或 Docker 继续坏着。
\end{itemize}

\begin{findingbox}{交付结果}
除现场修复外，仓库中的部署脚本也已经加固：\filepath{scripts/deploy.sh} 与 \filepath{scripts/deploy.ps1} 新增了跨平台的 \texttt{restart} / \texttt{probe-only} 模式，能自动打印 compose 状态、容器 health、端口映射、宿主机健康探测、容器内健康探测，并在失败时给出最近日志；对应 bash 路径已有单元测试覆盖。
\end{findingbox}

\section{研究范围与方法}

\begin{methodbox}{研究方法}
本报告只使用本地一手证据或当前工作树中的实现文件，不依赖事后转述。核心证据分为四类：
\begin{itemize}
\item 运行时配置：对比故障期间保留下来的 compose 备份与修复后的 live compose。
\item 运行时观测：\texttt{docker inspect}、\texttt{docker ps}、CCH 健康检查、host/container 网络探测。
\item 应用日志：验证 provider 选择、是否走代理、以及成功请求是否已回到 direct path。
\item 仓库实现：核对 deploy 脚本新增的 restart/probe 能力，以及 README 和单测是否同步。
\end{itemize}
\end{methodbox}

\subsection{证据边界}

本报告区分两类结论：

\begin{itemize}
\item \textbf{强证据结论}：能直接落到备份文件、当前容器元数据、当前命令输出或 repo 文件。
\item \textbf{次强结论}：例如 Docker NAT / \texttt{DOCKER} 链不一致，这在真实修复过程中确实被观察到，但当前没有像坏 \texttt{dns:} 一样保留下来的同等级原始证据。因此它在本报告中会被列为 ``可能放大故障的运行时副作用''，而不是最强根因。
\end{itemize}

\subsection{时间表示}

容器日志使用 UTC 时间戳，例如 \texttt{2026-03-25T12:29:18Z}；本地命令与 Docker service 观察使用 CST。为了避免 ``刚才''、``后来'' 这种模糊表述，正文尽量使用绝对日期和时间。

\section{这次 Bug 到底坏在哪里}

\subsection{先排除一个常见误判：provider 没有被选错}

如果问题真的出在 provider 配置本身，最典型的症状会是：

\begin{itemize}
\item provider selector 选不到 \texttt{RC/Hiyo}；
\item 选到了 provider，但 endpoint 变成了错误地址；
\item 请求还没到网络层就因为参数或 URL 拼接错误失败。
\end{itemize}

这次并不是这样。日志证据明确表明，在 \texttt{2026-03-25T12:29:18Z}，CCH 仍然做出了如下选择：

\begin{lstlisting}[caption={应用日志显示 provider 选择仍然是 RC，并尝试命中 right.codes}]
{"time":"2026-03-25T12:29:18.455Z","selected":{"name":"RC","id":1,"type":"codex"},"msg":"ProviderSelector: Selection decision"}
{"time":"2026-03-25T12:29:18.473Z","selectedEndpoints":[{"baseUrl":"https://right.codes/codex/"}],"msg":"ProxyForwarder: Trying provider"}
{"time":"2026-03-25T12:29:19.964Z","statusCode":200,"msg":"[ResponseHandler] Streaming request finalized as success"}
\end{lstlisting}

\begin{evidencebox}{判定}
因此，``provider 被改爆'' 这个表述太宽了。更精确的说法是：\textbf{provider 配置层仍能正确选到 RC，真正出问题的是容器访问公网 endpoint 时依赖的解析与网络路径。}
\end{evidencebox}

\subsection{最强根因：runtime compose 中的坏 DNS 覆盖}

保留下来的备份文件 \filepath{docker-compose.archbox.yaml.bak-20260325-1912} 给出了最直接证据。对应片段如下：

\begin{lstlisting}[caption={故障期间 runtime compose 里的坏 DNS 覆盖}]
ports:
  - "${APP_PORT:-23000}:${APP_PORT:-23000}"
dns:
  - 119.29.29.29
  - 223.5.5.5
extra_hosts:
  - "host.docker.internal:host-gateway"
\end{lstlisting}

这意味着 app 容器不再按 Docker 默认方式使用宿主机提供的解析链，而是被强制指向两台公网 DNS。对一个需要稳定出公网、并且还要与宿主机代理桥接共存的容器来说，这是高风险改动，原因至少有三层：

\begin{itemize}
\item 它绕过了 Docker 在当前主机上的默认解析链，增加了环境耦合。
\item 当网络策略、上游 DNS 路由、宿主机代理或 nftables/iptables 状态发生变化时，容器比宿主机更容易出现 ``宿主机能上网、容器却不稳定'' 的分裂。
\item 这类改动不会改变 provider 名称或 URL，所以表面看起来像 ``某些 provider 失灵''，但根本层级实际上更低。
\end{itemize}

\subsection{修复后的 live 容器已经不再携带该 DNS 覆盖}

当前 live 容器的 \texttt{docker inspect} 结果显示：

\begin{lstlisting}[caption={当前 live app 容器不再有 Docker 级别 Dns 覆盖}]
null ["host.docker.internal:100.121.76.120"] {"com.docker.compose.project.config_files":"/home/fanghaotian/Applications/claude-code-hub/docker-compose.archbox.yaml", ...}
\end{lstlisting}

这里第一个字段就是 \texttt{HostConfig.Dns}，其值为 \texttt{null}。这说明当前 live CCH 容器已不再带着那段坏配置运行。同时，当前运行目录下的两个 compose 文件都已经清理掉了 \texttt{dns:} 块。

\begin{evidencebox}{修复边界}
修复不是 ``把 provider 改回去''，而是 \textbf{把运行时容器恢复到不带坏 DNS 覆盖的状态}。这也是为什么部署脚本后续需要增加 runtime probe，而不是只围绕 provider 表做修修补补。
\end{evidencebox}

\section{为什么 Tailscale Serve 正常，不代表容器能出公网}

这是这次讨论里最容易混淆的一点。下图把两条链路拆开了：

\begin{figure}[H]
\centering
\begin{tikzpicture}[
    node distance=12mm and 10mm,
    box/.style={draw, rounded corners, align=center, minimum width=2.8cm, minimum height=1cm, fill=blue!5},
    warnbox/.style={draw, rounded corners, align=center, minimum width=3.1cm, minimum height=1cm, fill=red!4},
    arrow/.style={-{Latex[length=3mm]}, thick}
]
\node[box] (client) {其他电脑\\tailnet 客户端};
\node[box, right=of client] (serve) {Tailscale Serve\\宿主机入口};
\node[box, right=of serve] (hostloop) {宿主机 loopback\\127.0.0.1:23110 / 23000};
\node[box, right=of hostloop] (service) {live-dashboard / CCH\\本地服务入口};

\node[warnbox, below=18mm of hostloop] (container) {app 容器\\主动出公网};
\node[warnbox, right=of container] (dns) {Docker DNS / TLS\\出站链路};
\node[warnbox, right=of dns] (endpoint) {right.codes /\\codex.hiyo.top};

\draw[arrow] (client) -- (serve);
\draw[arrow] (serve) -- (hostloop);
\draw[arrow] (hostloop) -- (service);

\draw[arrow] (container) -- (dns);
\draw[arrow] (dns) -- (endpoint);

\node[align=left, above=4mm of serve, text width=6.6cm] {\small 上面这条链路证明的是 ``别人能不能打进来''。};
\node[align=left, below=4mm of dns, text width=6.6cm] {\small 下面这条链路证明的是 ``容器自己能不能解析并访问外部 endpoint''。};
\end{tikzpicture}
\caption{Tailscale Serve 的 ingress 路径与容器 egress 路径不是同一件事。}
\end{figure}

因此，``我看 Tailscale Serve 访问好像挺正常'' 这句话只证明：

\begin{itemize}
\item 宿主机上的 Serve 路由还在；
\item 宿主机 loopback 端口能接到请求；
\item 被转发到的本地服务入口没有完全死掉。
\end{itemize}

它并不能证明：

\begin{itemize}
\item app 容器的 DNS 解析链健康；
\item 容器出站 TCP/TLS 正常；
\item 某个外部 provider endpoint 的 TLS 行为正常。
\end{itemize}

\section{当前状态的精确结论}

\subsection{\texttt{right.codes} 现在对 host 和容器都正常}

当前命令输出显示：

\begin{itemize}
\item host 上执行 \texttt{curl -I -m 10 https://right.codes/} 返回 \texttt{HTTP/2 200}；
\item live app 容器里执行 \texttt{curl -I -m 10 --http1.1 https://right.codes/} 返回 \texttt{HTTP/1.1 200 OK}；
\item CCH 本地健康检查 \texttt{/api/actions/health} 返回 \texttt{\{"status":"ok"\}}；
\item app 容器处于 \texttt{healthy} 状态。
\end{itemize}

这足以支持一个保守但明确的判断：\textbf{对 RC 这条路径，当前 live CCH 的容器 egress 已经恢复到可用状态。}

\subsection{\texttt{codex.hiyo.top} 的现状是 endpoint 级残余，不是同一根因}

当前 host 和容器里都能解析出 \texttt{codex.hiyo.top}，但二者都会得到同样的 TLS EOF：

\begin{lstlisting}[caption={当前 host 与容器对 codex.hiyo.top 的共同症状}]
curl: (35) TLS connect error: error:0A000126:SSL routines::unexpected eof while reading
\end{lstlisting}

这里最重要的不是 ``它还失败''，而是 \textbf{host 和 container 失败得一样}。这会显著改变定位策略：

\begin{itemize}
\item 如果 host 能成、container 不能成，优先怀疑容器 DNS / egress；
\item 如果 host 和 container 都同样失败，优先怀疑 endpoint 自身、TLS 协商、站点策略或外部网络条件。
\end{itemize}

\begin{riskbox}{不要把两个问题混成一个}
``这次容器 DNS 覆盖导致 RC/Hiyo 无法直连'' 是主事故；``当前 \texttt{codex.hiyo.top} 在 host/container 同时 TLS EOF'' 是另一个更像 endpoint 侧的不稳定症状。后者不能反证前者没有发生，也不应该被误写成 ``CCH 现在还没修好 DNS''。
\end{riskbox}

\subsection{关于 Docker NAT / \texttt{DOCKER} 链}

真实修复过程中，Docker NAT 状态的不一致确实让某次重启操作变得更糟，并促成了 daemon restart。但是，与坏 DNS 覆盖不同，本报告没有保留到同等级的一手证据片段。因此本报告只做谨慎结论：

\begin{itemize}
\item 它很可能是 \textbf{放大故障和增加恢复成本} 的副问题；
\item 但最强、最可复核的根因证据，仍然是运行时 compose 的坏 \texttt{dns:} 覆盖。
\end{itemize}

\section{部署脚本加固了什么}

\subsection{bash 与 PowerShell 现在都支持 runtime 模式}

这次对部署脚本的目标不是 ``让一键安装更漂亮''，而是把 \textbf{已有部署的运行态修复} 做成可重复动作。对应实现位置：

\begin{itemize}
\item bash 参数与入口：\filepath{scripts/deploy.sh:66-94,135-182,1280-1282}
\item bash runtime probe / restart 核心：\filepath{scripts/deploy.sh:500-848}
\item bash 普通部署路径复用 probes：\filepath{scripts/deploy.sh:1123-1145,1305-1318}
\item PowerShell 参数与入口：\filepath{scripts/deploy.ps1:80-108,115-169,1334-1339}
\item PowerShell runtime probe / restart 核心：\filepath{scripts/deploy.ps1:400-860}
\item PowerShell 普通部署路径复用 probes：\filepath{scripts/deploy.ps1:1164-1198,1363-1380}
\end{itemize}

\subsection{新增能力总表}

\begin{longtable}{p{3.1cm}p{5.1cm}p{5.5cm}}
\toprule
能力 & 现在做法 & 直接收益 \\
\midrule
\endfirsthead
\toprule
能力 & 现在做法 & 直接收益 \\
\midrule
\endhead
已有部署重启 & 新增 \texttt{--restart} / \texttt{-Restart}，优先 \texttt{docker compose restart}，失败时回退到 \texttt{up -d} & 把 ``先试 restart，不行再 up -d'' 变成脚本内建逻辑 \\
纯探测模式 & 新增 \texttt{--probe-only} / \texttt{-ProbeOnly} & 无需改容器即可先看现场状态，减少误操作 \\
compose/env 显式定位 & 支持 \texttt{--compose-file}、\texttt{--env-file}、\texttt{--app-service} & 适配 \texttt{archbox} 这类 runtime compose 不在仓库内的部署 \\
宿主机探测 & 探测 \texttt{http://127.0.0.1:\{APP\_PORT\}/api/actions/health} & 直接判断本机入口是否可用 \\
容器内探测 & 用 \texttt{docker exec} 在容器内打同一 health 路径 & 区分 ``宿主机端口坏'' 与 ``容器内部其实是活的'' \\
运行态摘要 & 打印 \texttt{docker compose ps}、每个 service 的 state / health / ports & 从 ``黑盒重启'' 升级到 ``带可见性的重启'' \\
失败日志回显 & 不 ready 时自动打印最近日志 tail & 避免每次故障都手敲多条命令 \\
Windows 对齐 & PowerShell 版跟 bash 版保持基本一致的参数语义 & 文档和操作手册可以跨平台复用 \\
\bottomrule
\end{longtable}

\subsection{README 和测试也同步了}

\begin{itemize}
\item README 在 \filepath{README.md:167-183} 新增了 ``已有部署的重启与探测'' 小节，直接给出 Linux/macOS 与 Windows 示例。
\item bash 侧新增单测 \filepath{tests/unit/scripts/deploy.test.ts:251-295}，覆盖两条最关键路径：
  \begin{itemize}
  \item \texttt{probe-only} 会使用自定义 compose / env 并完成 host/container 探测；
  \item \texttt{restart} 在 \texttt{docker compose restart} 失败时会回退到 \texttt{up -d}。
  \end{itemize}
\end{itemize}

\begin{riskbox}{验证边界}
bash 版本已通过 \texttt{bash -n} 和 Vitest 单测。PowerShell 版本因为当前工作机没有 \texttt{pwsh}，只能做到静态人工审阅，尚未本机执行验证。这一点应该在合并或交付时明确写给维护者，而不是默认当成已跑通。
\end{riskbox}

\section{以后怎么修}

\subsection{推荐的排障顺序}

以后遇到 ``某个 provider 直连突然不行''，建议把排障顺序固定成下面这张表，而不是立刻改 provider 配置：

\begin{longtable}{p{4.1cm}p{3.3cm}p{4.1cm}p{2.8cm}}
\toprule
观察到的现象 & host 探测 & container 探测 & 优先动作 \\
\midrule
\endfirsthead
\toprule
观察到的现象 & host 探测 & container 探测 & 优先动作 \\
\midrule
\endhead
Tailscale Serve 正常，但 provider 请求失败 & 不足以判断 & 不足以判断 & 先跑 runtime probe，不要把 ingress 当 egress \\
host 能 \texttt{curl}，container 不能 \texttt{curl} & 成功 & 失败 & 优先看 container DNS、compose 覆盖项、Docker egress \\
host 和 container 都失败，而且报错一致 & 失败 & 失败 & 优先怀疑 endpoint 自身、TLS 协商、上游网络条件 \\
应用日志仍选中正确 provider，但报 \texttt{EAI\_AGAIN} & 可能成功也可能失败 & 可能失败 & 优先看 DNS / 容器网络层，不要先改 provider URL \\
日志出现 \texttt{Using proxy} & 不相关 & 不相关 & 说明请求没有直连；先确认是否还残留 \texttt{proxy\_url} 或代理桥接配置 \\
\bottomrule
\end{longtable}

\subsection{推荐命令}

\begin{lstlisting}[language=bash, caption={Linux/macOS 上建议的第一轮排障命令}]
./scripts/deploy.sh --probe-only -d /home/fanghaotian/Applications/claude-code-hub
curl -I -m 10 https://right.codes/
docker exec claude-code-hub-app-s6gz sh -lc 'curl -I -m 10 --http1.1 https://right.codes/'
\end{lstlisting}

\begin{lstlisting}[language={}, caption={Windows 上建议的第一轮排障命令}]
.\scripts\deploy.ps1 -ProbeOnly -DeployDir C:\ProgramData\claude-code-hub
\end{lstlisting}

\subsection{真正要避免的不是 ``重启''，而是 ``盲重启''}

这次事故暴露出来的最大流程问题，不是没人会重启，而是重启前后缺少统一的现场摘要：

\begin{itemize}
\item 不知道当前到底用了哪一个 compose 文件；
\item 不知道 app service 的健康状态和端口映射；
\item 不知道宿主机入口活不活；
\item 不知道容器内部 health 活不活；
\item 不知道失败后最近的日志应该先看哪里。
\end{itemize}

脚本加固的价值恰恰在这里：\textbf{先把现场抽象成固定的 probe，再决定是否重启}。

\section{结论与后续问题}

\subsection{已被强证据支持的结论}

\begin{itemize}
\item 本次主事故的最强根因是 runtime compose 中的坏 \texttt{dns:} 覆盖，而不是 provider base URL 被改错。
\item Tailscale Serve 正常不能推出容器可以出公网，因为它验证的是 ingress，不是 egress。
\item 当前 live CCH 容器已经不再带 Docker 级别 \texttt{Dns} 覆盖，CCH 健康检查正常，\texttt{right.codes} 从 host 和容器都可访问。
\item bash 与 PowerShell 部署脚本已经新增 runtime probe / restart 能力，README 和 bash 单测也已同步。
\end{itemize}

\subsection{大概率成立，但证据弱于主结论的判断}

\begin{itemize}
\item Docker NAT / \texttt{DOCKER} 链的不一致在真实修复过程中确实增加了恢复复杂度，但它更像次级放大因素，而不是主事故的最佳解释。
\end{itemize}

\subsection{仍待继续观察的问题}

\begin{itemize}
\item \texttt{codex.hiyo.top} 当前在 host 和 container 都报 TLS EOF，这说明 Hiyo 端点现在还有一个与本次 DNS 事故不同的残余问题；如果业务仍依赖它，应该按 ``endpoint/TLS 问题'' 单独开案处理。
\item \texttt{scripts/deploy.ps1} 需要在有 \texttt{pwsh} 的 Windows 机器上做一次真实运行验证，当前只能说实现已对齐、尚未本机执行。
\end{itemize}

\appendix

\section{来源清单}

\begin{longtable}{p{1.0cm}p{1.8cm}p{4.2cm}p{7.3cm}}
\toprule
ID & 模态 & 标题 / 标签 & 路径或说明 \\
\midrule
\endfirsthead
\toprule
ID & 模态 & 标题 / 标签 & 路径或说明 \\
\midrule
\endhead
S1 & 配置备份 & 坏 DNS compose 备份 & \filepath{/home/fanghaotian/Applications/claude-code-hub/docker-compose.archbox.yaml.bak-20260325-1912} \\
S2 & 配置 & 当前 runtime compose 状态 & \filepath{/home/fanghaotian/Applications/claude-code-hub/docker-compose.archbox.yaml}；\filepath{/home/fanghaotian/Applications/claude-code-hub/docker-compose.yaml} \\
S3 & 运行时 & live 容器元数据 & \texttt{docker inspect claude-code-hub-app-s6gz --format ...} \\
S4 & 运行时 & 健康检查与容器状态 & \texttt{docker compose ps}、\texttt{/api/actions/health}、\texttt{docker ps --format} \\
S5 & 日志 & RC 请求日志 & \texttt{docker logs claude-code-hub-app-s6gz} \\
S6 & 网络探测 & host/container 连通性核验 & \texttt{curl} 与 \texttt{docker exec ... curl} \\
S7 & 代码 & bash deploy 脚本 & \filepath{scripts/deploy.sh} \\
S8 & 代码 & PowerShell deploy 脚本 & \filepath{scripts/deploy.ps1} \\
S9 & 代码 & bash 单测 & \filepath{tests/unit/scripts/deploy.test.ts} \\
S10 & 文档 & README restart/probe 入口 & \filepath{README.md} \\
S11 & 备注 & 原始命令输出汇总 & \filepath{docs/reports/2026-03-25-cch-network-postmortem/incident-evidence.md} \\
\bottomrule
\end{longtable}

\end{document}
